% Portada i índex del recull d'enunciats de càlcul.
% Els compila scripts/munta_recull_pdf.py amb XeLaTeX i els posa davant del PDF del recull.
% L'script genera index_recull.tex (entrades de l'índex amb el número de pàgina) i
% parametres_recull.tex (text d'edició i mida de l'índex). Sense aquests fitxers, només surt la portada.
% Fa servir l'estil dels apunts (apunts.sty) perquè coincideixi amb la portada dels PDF dels temes.
\documentclass[11pt,a4paper,oneside]{book}
\usepackage{apunts}
\hypersetup{pdftitle={Recull d'enunciats de càlcul},pdfauthor={Albert Cid},pdflang={ca}}
\pagestyle{empty}
\usepackage{multicol}
\setlength{\columnsep}{2.2em}

\newcommand{\edicio}{}
\newcommand{\midaindex}{\normalsize}
\InputIfFileExists{parametres_recull.tex}{}{}

% \entradaindex{nivell}{títol}{pàgina}: nivell 2 (bloc), 3 (apartat) o 4 (subapartat)
\newcommand{\puntsindex}{\nobreak\leaders\hbox to 0.7em{\hss.\hss}\hfill\nobreak}
\newcommand{\paginaindex}[1]{\makebox[1.8em][r]{#1}}
\newcommand{\entradaindex}[3]{%
  \ifnum#1=2
    \par\vspace{0.7em}\noindent\hangindent=1em{\bfseries\color{tema}#2\puntsindex\paginaindex{#3}}\par\vspace{0.2em}%
  \else\ifnum#1=3
    \par\noindent\hangindent=1.8em\hspace*{0.8em}#2\puntsindex\paginaindex{#3}\par
  \else
    \par\noindent\hangindent=2.8em\hspace*{1.8em}{\color{black!70}#2}\puntsindex{\color{black!70}\paginaindex{#3}}\par
  \fi\fi}

\begin{document}
\begin{titlepage}
  \sffamily\color{tema}
  \vspace*{0.16\textheight}
  {\Large Matemàtiques Aplicades a les Ciències Socials II\par}
  \vspace{0.4em}
  {\large\color{black!60}2n de Batxillerat\par}
  \vspace{3.2em}
  {\color{tema}\rule{\linewidth}{1.2pt}\par}
  \vspace{1.4em}
  {\Large\MakeUppercase{Càlcul}\par}
  \vspace{0.5em}
  {\fontsize{38}{44}\selectfont\bfseries Recull d'enunciats\par}
  \vspace{1.4em}
  {\color{tema}\rule{\linewidth}{1.2pt}\par}
  \vspace{2.2em}
  \autoria
  \vspace{3em}
  {\large\color{black!80}
  \begin{tabular}{@{}l@{\hspace{1em}}l@{}}
    \color{tema}Tema 1 & Repàs de límits i derivades\\[0.5em]
    \color{tema}Tema 2 & Aplicacions de les derivades\\[0.5em]
    \color{tema}PAU & Problemes de càlcul
  \end{tabular}\par}
  \vfill
  {\normalsize\color{black!55}Recull d'exercicis\ifx\edicio\empty\else\ · \edicio\fi. Les solucions pas a pas són al web, a
  \href{https://albertcid-pa.github.io/Matematiques2BAT-MACS/calcul/per-practicar/exercicis_calcul/}{Càlcul\,›\,Per practicar}.\par}
  \vspace{1.2em}
  {\color{black!25}\rule{\linewidth}{0.4pt}\par}
  \vspace{0.8em}
  \llicencia
\end{titlepage}

\IfFileExists{index_recull.tex}{%
  \newpage
  {\sffamily\midaindex\setlength{\parskip}{0pt}\linespread{1.25}\selectfont
   {\color{tema}\fontsize{24}{28}\selectfont\bfseries Índex\par}
   \vspace{0.3em}
   {\color{tema}\rule{\linewidth}{1.2pt}\par}
   \vspace{0.2em}
   {\color{black!55}\small Els números de pàgina són els del peu de cada pàgina d'exercicis.\par}
   \vspace{0.4em}
   \begin{multicols}{2}\raggedright\input{index_recull.tex}\end{multicols}\par}
}{}
\end{document}
